% sections/08_exploratory_boxplots.tex — V_B1, V_B2, V_B3  [NEW in v3]

\section{Exploratory Distribution Analysis (Box Plots)}

Prior to hypothesis testing, three 5-point summary box plots were generated to
understand the distributional properties of key quantitative variables and
validate the $2\times$IQR outlier removal strategy. These charts informed data
cleaning decisions and reveal structural characteristics that qualify all
downstream analysis.

%% ── V_B1 ─────────────────────────────────────────────────────
\subsection{V\_B1: Accident Distance Distribution by State}

A side-by-side box plot displays the distribution of \texttt{Distance(mi)} ---
the road length affected by the accident --- categorized by all 49 states.
Grey box plots show post-cleaning values; blue dotted whisker extensions show
pre-cleaning outliers, making the effect of $2\times$IQR removal explicit.

% TO ADD FIGURE: save as figures/vb1_distance_by_state.png
% \figchart[0.98]{vb1_distance_by_state}{Accident distance by state pre- and post-cleaning}{fig:vb1}
\begin{figure}[H]
  \centering
  \figchart[0.98]{vb1_distance_by_state}{Accident distance distribution by state post $2\times$IQR cleaning}{fig:vb1}
  \caption{}
  \label{fig:vb1}
\end{figure}

\idiomtable
  {1D Line (box: Q1, median, Q3) + 0D Points (whiskers, outliers)}
  {Vertical Length (ordered, aligned)}
  {Analyse distribution spread; compare outlier removal effect across states}
  {49 states --- manageable for a side-by-side layout with small font labels}

\insightbox{Most states show a median accident distance of 0.5--1.5 miles (the
IQR box), confirming the dataset predominantly captures localised road blockages.
Georgia (GA) is the most notable outlier with max whisker $\approx$7\,mi;
Montana (MT) and Wyoming (WY) show wide IQR spreads reflecting sparse rural
highway networks where accidents affect longer road stretches. Wyoming's
pre-clean maximum approaches 8.5 miles --- physically implausible for a single
accident. The $2\times$IQR cleaning effectively separates sensor errors from
genuine high-distance incidents, directly justifying the outlier removal step.}

%% ── V_B2 ─────────────────────────────────────────────────────
\subsection{V\_B2: Accident Duration Distribution by State}

A side-by-side box plot shows the distribution of \texttt{Duration}
(time from accident start to traffic clearance) per state. This is the most
direct continuous proxy for the severity score --- longer duration implies a
higher severity classification.

% TO ADD FIGURE: save as figures/vb2_duration_by_state.png
% \figchart[0.98]{vb2_duration_by_state}{Accident duration by state pre- and post-cleaning}{fig:vb2}
\begin{figure}[H]
  \centering
    \figchart[0.98]{vb2_duration_by_state}{Accident duration distribution by state post-$2\times$IQR cleaning}{fig:vb2}
  \label{fig:vb2}
\end{figure}

\idiomtable
  {1D Line (box) + 0D Points (whiskers)}
  {Vertical Length (ordered, aligned)}
  {Analyse distributional spread; validate outlier removal on Duration}
  {49 states --- same layout as V\_B1 for direct comparison}

\insightbox{The typical accident resolves in 75--120 minutes (median per state),
with an IQR of $\approx$55--170 minutes. This is the continuous signal underlying
the ordinal Severity score --- Severity~4 records occupy the upper tail.
Maine's pre-clean maximum of $\approx$950 minutes ($\approx$16 hours) is almost
certainly a timestamp entry error: a 16-hour road closure would be a major media
event. Oregon's wide IQR may reflect genuine longer clearance times on rural
highways where emergency response is slower --- exactly the mechanism supporting H4.
The $2\times$IQR cleaning removed the most extreme implausible values while
retaining genuine long-duration incidents. Crucially, the consistent cross-state
box shape (similar Q1--Q3 widths) confirms the dataset is structurally sound
after cleaning.}

%% ── V_B3 ─────────────────────────────────────────────────────
\subsection{V\_B3: Accident Count Distribution by Year (2016--2022)}

A box plot with overlaid state-level dots shows the per-state accident count
distribution for each year 2016--2022 over weeks. The COVID year 2020 is highlighted in
yellow. This chart simultaneously validates dataset growth (Fault~7), exposes
geographic inequality (Fault~2), and contextualises the COVID anomaly.

% TO ADD FIGURE: save as figures/vb3_count_by_year.png
% \figchart{vb3_count_by_year}{Per-state accident count by year with COVID highlight}{fig:vb3}
\begin{figure}[H]
  \centering
  \figchart{vb3_count_by_year}{Per-state accident count distribution by year (2016--2022) over weeks}{fig:vb3}
  \label{fig:vb3}
\end{figure}

\idiomtable
  {1D Line (box) + 0D Points (state-level dots)}
  {Vertical Length (ordered, aligned) + Colour Hue (yellow = COVID 2020)}
  {Identify trend in dataset growth; expose geographic inequality}
  {7 years $\times$ 52 week dots --- moderate density, overplotting managed by jitter}

\insightbox{Three simultaneous findings emerge. First, the median per-state
accident count grows steadily from $\approx$8K (2016) to $\approx$26K (2022),
confirming the dataset expansion noted in Fault~7 --- more states being covered
each year, not just more accidents. Second, the box grows taller each year (wider
IQR), revealing increasing geographic inequality: some states (CA, FL, TX) grow
much faster, widening the distribution.
Excluding 2023 is clearly
justified: a partial-year box cannot be meaningfully compared to these full
annual distributions.}
